\begin{lemma}[Smooth1Properties-above]
\uses{Smooth1}
\label{Smooth1Properties_above}
\lean{Smooth1Properties_above}
% at /home/runner/work/trough/trough/library/PrimeNumberTheoremAnd/MellinCalculus.lean:1011.0-1074.54
\leanok
Fix $0<\epsilon<1$. There is an absolute constant $c>0$ so that:
  if $x\geq (1+c\epsilon)$, then
  $$\widetilde{1_{\epsilon}}(x) = 0.$$
\end{lemma}
\begin{proof}
\uses{DeltaSpike,MellinConvolution}
\leanok
Again the Mellin convolution is
  $$\int_0^1 \nu_\epsilon(x/y)\frac{dy}{y},$$
  but now if $x > 2^\epsilon$, then the support of $\nu_\epsilon$ is disjoint
  from the region of integration, and hence the integral is zero.
  We choose
  $$
    c:=2\log 2
    .
  $$
  By Lemma \ref{Smooth1Properties_estimate},
  $$
    c > 2\frac{1-2^{-\epsilon}}\epsilon > 2^\epsilon\frac{1-2^{-\epsilon}}\epsilon
    =
    \frac{2^\epsilon-1}\epsilon,
  $$
  so
  $$
    1+c\epsilon > 2^\epsilon.
  $$
\end{proof}
